\begin{tikzpicture}[
    >=Latex,
    scale=1.4,
    every node/.style={font=\small},
]
    % --- Buildplate ---
    \fill[gray!15] (-2.5, -0.4) rectangle (3.0, 0);
    \draw[gray!60] (-2.5, -0.4) rectangle (3.0, 0);

    % --- Bauteil mit unregelmäßiger Form (Pilz-Topf-artig) ---
    % Linker Teil: glatte Wölbung
    % Rechter Teil: Schulter mit Lücke darüber
    \fill[blue!10]
        (-2.0, 0) -- (-2.0, 0.4)
        .. controls (-1.8, 1.5) and (-0.8, 2.1) ..
        (0.5, 2.0) --
        (1.0, 1.6) --
        (1.0, 0.5) --
        (1.5, 0.5) --
        (1.5, 0) -- cycle;
    \draw[thick, blue!60!black]
        (-2.0, 0) -- (-2.0, 0.4)
        .. controls (-1.8, 1.5) and (-0.8, 2.1) ..
        (0.5, 2.0) --
        (1.0, 1.6) --
        (1.0, 0.5) --
        (1.5, 0.5) --
        (1.5, 0);
    \node[blue!60!black, font=\scriptsize, anchor=west]
        at (-2.0, -0.6) {Bauteil};

    % --- Mittelpunkt c ---
    \fill[red!80!black] (-0.3, 0) circle (2pt);
    \node[red!80!black, anchor=north east, font=\scriptsize]
        at (-0.3, -0.05) {$\vec{c}$};

    % --- Sphärenschale S_i ---
    \draw[thick, orange!85!black]
        ($(-0.3, 0) + (0:1.7)$) arc (0:180:1.7);
    \node[orange!85!black, font=\scriptsize, anchor=west]
        at ($(-0.3, 0) + (0:1.75)$) {$S_i$};

    % =========================================================
    % Punkt A (druckbar, im Bauteil)
    % auf der Sphäre bei Winkel ~110°
    % =========================================================
    \coordinate (A) at ($(-0.3, 0) + (115:1.7)$);
    \fill[green!50!black] (A) circle (2.5pt);
    \node[green!50!black, anchor=south east, font=\scriptsize]
        at ($(A) + (-0.05, 0.05)$) {$\vec{p}_A$};

    % Radiale Normale n_A (nach außen)
    \draw[->, thick, green!50!black]
        (A) -- ($(A) + (115:0.45)$);
    \node[green!50!black, font=\scriptsize, anchor=south east]
        at ($(A) + (115:0.5)$) {$\hat{n}_A$};

    % Test-Punkt p_A - delta * n_A (nach innen)
    \fill[green!50!black] ($(A) + (115:-0.3)$) circle (1.5pt);
    \draw[->, thin, green!50!black, dashed]
        (A) -- ($(A) + (115:-0.3)$);
    \node[green!50!black, font=\tiny, anchor=west,
            align=left, fill=white, inner sep=1pt]
        at ($(A) + (115:-0.45)$)
        {im Bauteil\\$\Rightarrow$ druckbar};

    % =========================================================
    % Punkt B (Überhang, außerhalb Bauteil)
    % auf der Sphäre über der Lücke (Schulter rechts)
    % =========================================================
    \coordinate (B) at ($(-0.3, 0) + (35:1.7)$);
    \fill[red!80!black] (B) circle (2.5pt);
    \node[red!80!black, anchor=south west, font=\scriptsize]
        at ($(B) + (0.05, 0.05)$) {$\vec{p}_B$};

    % Radiale Normale n_B (nach außen)
    \draw[->, thick, red!80!black]
        (B) -- ($(B) + (35:0.45)$);
    \node[red!80!black, font=\scriptsize, anchor=south west]
        at ($(B) + (35:0.5)$) {$\hat{n}_B$};

    % Test-Punkt p_B - delta * n_B (nach innen, AUSSERHALB Bauteil)
    \fill[red!80!black] ($(B) + (35:-0.3)$) circle (1.5pt);
    \draw[->, thin, red!80!black, dashed]
        (B) -- ($(B) + (35:-0.3)$);
    \node[red!80!black, font=\tiny, anchor=north west,
            align=left, fill=white, inner sep=1pt]
        at ($(B) + (35:-0.4) + (0.05, -0.1)$)
        {außerhalb\\$\Rightarrow$ \glsentrytext{overhang}};

\end{tikzpicture}